\documentclass[reqno]{exam}

\usepackage{amssymb, amsmath, amsthm, stmaryrd}
\usepackage{fullpage, parskip}
\usepackage{graphicx}
\usepackage{enumitem}

% Improved spacing
\usepackage[top=1in, bottom=0.8in, left=1in, right=1in]{geometry}


\title{\normalsize{Quiz 9, ICM, Spring 2026 \hspace{1ex} JHED: \underline{\hspace{20ex}} \hspace{2ex} Name: \underline{\hspace{40ex}}}}
\author{}
\date{}

% Custom spacing for better page utilization
\renewcommand{\baselinestretch}{0.95}

\begin{document}
\maketitle
\vspace{-2cm} % Reduce space after title

\begin{questions}
  \question
  Consider the ordinary differential equation
  \begin{align*}
    u'(t) = \dfrac{1}{u - 2}.
  \end{align*}
  \begin{parts}
    \part[4] Is the ODE guaranteed to have a unique solution on the time interval $0 \le t \le 1$ for all initial conditions? Justify your answer.
    \part[4] Verify that the function $u(t) = 2 - \sqrt{1 + 2t}$ is a solution to the ODE with the initial condition $u(0) = 1$.
    \part[2] Explain how part (b) does not contradict the conclusion from part (a).
  \end{parts}

  \begin{solution}
    The ODE $u' = f(t, u)$ where $f(t, u) = \frac{1}{u-2}$ does not satisfy the theorem. The partial derivative
    $$\frac{\partial f}{\partial u} = -\frac{1}{(u-2)^2}$$
    is not bounded on $t \in [0, 1]$ for all $u$. Specifically, as $u \to 2$, $|\frac{\partial f}{\partial u}| \to \infty$, so a unique solution is not guaranteed.

    Given $u(t) = 2 - \sqrt{1 + 2t}$ and $u(0) = 2 - \sqrt{1} = 1$:
    \begin{align*}
      u'(t)         & = -\frac{1}{2\sqrt{1+2t}}(2) = -\frac{1}{\sqrt{1+2t}}      \\
      \frac{1}{u-2} & = \frac{1}{(2 - \sqrt{1+2t}) - 2} = \frac{1}{-\sqrt{1+2t}}
    \end{align*}
    Since LHS = RHS, the function is a solution.

    Part (b) does not contradict (a) because the theorem provides a sufficient condition, not a necessary one. The failure of the Lipschitz bound means uniqueness is not guaranteed by this specific theorem, but it does not preclude a solution from existing.

  \end{solution}

  \newpage
  \question[10]
  Consider the initial differential equation
  \begin{align*}
    u'(t) = \dfrac{1}{u - 2}
  \end{align*}\
  with initial condition $u(0) = 1$.

  Run the forward Euler method with step size $h = 0.5$ to compute approximations to $u(0.5)$ and $u(1)$. Show your work.

  \begin{solution}

    Given $h = 0.5$, $t_0 = 0$, $u_0 = 1$, and $f(t, u) = \frac{1}{u-2}$:

    \begin{align*}
      u_1 & = u_0 + h \cdot f(t_0, u_0)                                                          \\
      u_1 & = 1 + 0.5 \left( \frac{1}{1-2} \right) = 1 - 0.5 = 0.5                               \\
      u_2 & = u_1 + h \cdot f(t_1, u_1)                                                          \\
      u_2 & = 0.5 + 0.5 \left( \frac{1}{0.5-2} \right) = 0.5 + 0.5 \left( -\frac{1}{1.5} \right) \\
      u_2 & = \frac{1}{2} - \frac{1}{3} = \frac{1}{6}
    \end{align*}

    Therefore, $u(0.5) \approx 0.5$ and $u(1.0) \approx 0.1667$.
  \end{solution}


\end{questions}
\end{document}